\question{}

Consider Poisson's equation
\[
\frac{\partial^2 u}{\partial x^2} + \frac{\partial^2 u}{\partial y^2} + \frac{\partial^2 u}{\partial z^2} = 1
\]
in a cube $\Omega = \{(x,y,z)\in\mathbb{R}^3: 0\leq (x,y,z) \leq L\}$.
The boundary conditions are zero Dirichlet conditions. Consider an axis-aligned grid with steps $\Delta_x = L/J_x$, $\Delta_y = L/J_y$, $\Delta_z = L/J_z$.

%\begin{enumerate}[label=2.\arabic*]
\begin{parts}
%[label=\Alph*]
\part Write a finite volume scheme for this problem based on the given grid. 
    \begin{solutionbox}{18cm}
        Let $U_{i,j,k}$ denote the approximate solution of the PDE at vertex $(i \Delta_x, j \Delta_y, k \Delta_z) \in \Omega$ for $i = 0,\dots,J_x$, $j = 0,\dots,J_y$, $k = 0,\dots,J_z$. 
The scheme is
% \begin{multline*}
% \frac{U_{i+1,j,k} - 2 U_{i,j,k} + U_{i-1,j,k}}{\Delta_x} \Delta_y \Delta_z
% + \frac{U_{i,j+1,k} - 2 U_{i,j,k} + U_{i,j-1,k}}{\Delta_y} \Delta_z \Delta_x
% \\ + \frac{U_{i,j,k+1} - 2 U_{i,j,k} + U_{i,j,k-1}}{\Delta_z} \Delta_x \Delta_y
%  = \Delta_x \Delta_y \Delta_z
% \end{multline*}
\begin{multline*}
\frac{U_{i+1,j,k} - U_{i,j,k}}{\Delta_x} \Delta_y \Delta_z + \frac{U_{i-1,j,k} - U_{i,j,k}}{\Delta_x} \Delta_y \Delta_z
\\ + \frac{U_{i,j+1,k} - U_{i,j,k}}{\Delta_y} \Delta_z \Delta_x + \frac{U_{i,j-1,k} - U_{i,j,k}}{\Delta_y} \Delta_z \Delta_x
\\ + \frac{U_{i,j,k+1} - U_{i,j,k}}{\Delta_z} \Delta_x \Delta_y + \frac{U_{i-1,j,k-1} - U_{i,j,k}}{\Delta_z} \Delta_x \Delta_y
\\ = \Delta_x \Delta_y \Delta_z
\end{multline*}
for $i = 1,\dots,J_x-1$, $j = 1,\dots,J_y-1$, $k = 1,\dots,J_z-1$, with $U_{i,j,k} = 0$ whenever $i=0$ or $i=J_x$ or $j=0$ or $j=J_y$ or $k=0$ or $k=J_z$. 

\textcolor{blue}{Observe that the scheme simplifies to 
\begin{multline*}
\frac{U_{i+1,j,k} - 2 U_{i,j,k} + U_{i-1,j,k}}{\Delta_x^2}
+ \frac{U_{i,j+1,k} - 2 U_{i,j,k} + U_{i,j-1,k}}{\Delta_y^2}
\\ + \frac{U_{i,j,k+1} - 2 U_{i,j,k} + U_{i-1,j,k-1}}{\Delta_z^2}
 = 1.
\end{multline*}
}
    \end{solutionbox}
\part Explain how you obtained the scheme. 
    \begin{solutionbox}{20cm}
        \textcolor{blue}{A justification can be deduced from the part on finite volume methods in the course notes, considering the specified PDE, domain, boundary conditions, and the vertices $(i \Delta_x, j \Delta_y, k \Delta_z) \in \Omega$ for $i = 0,\dots,J_x$, $j = 0,\dots,J_y$, $k = 0,\dots,J_z$. Observe that with these vertices, the Voronoi cells are axis-aligned rectangular parallelepipeds: $V_{i,j,k} = \{(x,y,z) \in \Omega: (i-\frac12)\Delta_x \leq x \leq (i+\frac12)\Delta_x, (j-\frac12)\Delta_y \leq y \leq (j+\frac12)\Delta_y, (k-\frac12)\Delta_z \leq z \leq (k+\frac12)\Delta_z\}$ for $i = 0,\dots,J_x$, $j = 0,\dots,J_y$, $k = 0,\dots,J_z$. In other words, $V_{i,j,k}$ is the axis-aligned rectangular parallelepiped centered at $(i \Delta_x, j \Delta_y, k \Delta_z)$ of size $\Delta_x \times \Delta_y \times \Delta_z$ and restricted to $\Omega$, this restriction only affecting the cells of the vertices that belong to the boundary $\partial \Omega$. Parallelepipeds have six faces, and each of the six terms on the left-hand side of the scheme equation given in answer 2(a) corresponds to the contribution of the boundary part between cell $V_{i,j,k}$ and one of the six neighboring cells. 
} 
    \end{solutionbox}
\end{parts}
